% RESEARCH_DRAFT_SHA256: f763bdf8813aae5fc3f59bc0e92e1d34395bb81d8033ddf80f2adbb3f46a0c00
% Layout and source-reference wrappers expanded for the web renderer.
\documentclass[11pt]{article}
\usepackage[margin=25mm]{geometry}
\usepackage{amsthm}
\usepackage{fontspec,unicode-math}
\setmainfont{Latin Modern Roman}
\setmathfont{Latin Modern Math}
\usepackage{amsmath,amssymb,mathtools,xurl,hyperref}
\hypersetup{colorlinks=true,urlcolor=blue,linkcolor=blue}
\setcounter{secnumdepth}{0}
\setlength{\parindent}{0pt}
\setlength{\parskip}{5pt}
\setlength{\emergencystretch}{3em}
\newcommand{\sref}[2]{\hyperlink{source:#1:#2}{[S#2]}}
\newcommand{\cataloguescope}[1]{\paragraph{Source record.}#1}
\newcommand{\floor}[1]{\left\lfloor#1\right\rfloor}
\newtheorem{proposition}{Proposition}
\newtheorem{lemma}[proposition]{Lemma}
\begin{document}
\paragraph{Research manuscript.} OpenAI \texttt{gpt-6-astra}, reasoning \texttt{ultra}; 11 October 2026.
\paragraph{Preserved source body.} The complete catalogue mathematical text and sources follow with layout and source links expanded, including source qualifications and its historical preparation notes. New research begins at the section headed Research attempt.

% UnsolvedMath ID 30002625; source catalogue code OWR-13105-002
\section{Hardness Amplification by Exponent Noise}\label{30002625}
{\small UnsolvedMath ID 30002625 · Computer Science}
\subsection{Definitions and mathematical statement}

\paragraph{Shared notation.}Standard number systems are $\mathbb N,\mathbb Z,\mathbb Q,\mathbb R,\mathbb C$. Specialized notation and index ranges follow the cited source and the definitions below.


An exponent-noise learning problem replaces a modular linear observation by a group element encoding that observation in its exponent. Its difficulty depends on the group, exponent modulus, error distribution, and access to group representations.
\paragraph{Record scope and qualifications.}
The catalogue title is retained as metadata. Named certificates, contracts, bounds or frameworks in that title are not asserted as results of the original source extract.
The exact LWEE distribution and comparison/reduction notion are those of the report.

Source-record limitation: The exponent-noise learning task and hardness-comparison notion are not defined in the supplied question.
\paragraph{Statement recorded in the supplied source.}
Does adding noise to the exponent in a group-theoretic learning problem amplify computational intractability, as proposed by Learning with Errors in the Exponent? \href{https://doi.org/10.4171/owr/2014/35}{[S2]}
\paragraph{Source status.}
The supplied catalogue labels this source record partially solved; that metadata was checked against the supplied snapshot on 2026-10-11. The label is preserved for provenance and is not an independent verification of present open status or mathematical correctness. The catalogue literature assessment is dated 2026-08-21.
\subsection{Short English statement}
Determine whether exponent noise amplifies the hardness of the cited group learning task.
\subsection{Sources}
\begin{itemize}
\item[S1] Ulam AI and the UnsolvedMath Contributors, UnsolvedMath: A Curated Collection of Open Mathematics Problems, record 30002625 (OWR-13105-002). Catalogue curation and metadata: CC BY 4.0. The source extract is retained; definitions, English summaries and qualifications are editorial additions.
\url{https://huggingface.co/datasets/ulamai/UnsolvedMath}.
\item[S2] Cryptography (2015), Oberwolfach Reports, DOI 10.4171/owr/2014/35, pp. 1933--1990
\url{https://doi.org/10.4171/owr/2014/35}.
\item[S3] Dagdelen--Gajek--Gopfert, Learning with Errors in the Exponent
\url{https://eprint.iacr.org/2014/826.pdf}.
\item[S4] published LWEE paper
\url{https://doi.org/10.1007/978-3-319-30840-1_5}.
\item Source-specific Oberwolfach Reports licences, including CC BY-SA 4.0 where applicable, are retained; see THIRD\_PARTY\_NOTICES.md.
\end{itemize}
\paragraph{Preparation and review.}
The original source and UnsolvedMath contributors retain attribution for the source extract. Codex (GPT-6; exact serving revision not exposed) assisted with definitions, English summaries, scope notes and formatting. No new human mathematical authorship or independently verified solution is asserted. Permanent source ID 30002625 and code OWR-13105-002 checked against the complete supplied registry and current repository. Distinct upstream records, including repeated questions and broad or damaged extracts, are retained for complete coverage of the source CS category. A repeated record is not a new independent conjecture. Complete mathematical-equivalence and current-literature review has not been performed. Preparation and final review follow contribution prompts C and D, with source gaps disclosed. Structural and rendering checks do not establish a new theorem or open status.
\label{end:30002625}

\section{Research attempt}
\subsection{Literature, source context and scope}
The full original catalogue body and source qualifications are preserved above. The relevant primary workshop passages were read in the cached publisher report. Historical claims and subsequent attributed results are distinguished from the mathematical work below. A restricted example, standard lemma or known theorem is not counted as new progress on the entire original target. Missing assumptions and repeated or rhetorical extracts are not silently turned into different solved problems.

\subsection{Noise contraction is not a computational hardness separation}
The actual learning-with-errors-in-the-exponent samples are $(g^a,g^{\langle a,s\rangle+e})$, with parameters and a nonuniform error law specified in the source. We test the proposed inference that adding group noise must by itself create a strictly harder problem. The following finite-group calculation is information theoretic and deliberately does not assert a computational reduction between the source problems.
\begin{proposition}
For distributions $P,Q,\nu$ on a finite additive group, convolution with common independent noise obeys
\[
\operatorname{TV}(P*\nu,Q*\nu)\le\operatorname{TV}(P,Q).
\]
A fixed deterministic noise preserves this distance exactly; uniform noise makes the two output distributions identical.
\end{proposition}
\begin{proof}
Using $\operatorname{TV}(P,Q)=\tfrac12\sum_y|P(y)-Q(y)|$, the triangle inequality gives
\[
\tfrac12\sum_y\left|\sum_z\nu(z)(P(y-z)-Q(y-z))\right|
\le\tfrac12\sum_z\nu(z)\sum_y|P(y-z)-Q(y-z)|
=\operatorname{TV}(P,Q).
\]
For a point-mass noise at $z_0$, translation simply permutes the summands, proving equality. If $\nu$ is uniform, $(P*\nu)(y)=|G|^{-1}\sum_xP(x)=|G|^{-1}$ for every $y$, independently of $P$, and likewise for $Q$.
\end{proof}
For a cyclic group of order $q$ with generator $g$, exponentiation is a bijection of the underlying finite sets, so the same statistical statements hold for $g^{X+e}$. If fresh $e$ is uniform modulo $q$, then $g^{\langle a,s\rangle+e}$ is uniform conditional on $a,s$. Even arbitrarily many such independently noised samples convey no information about $s$ through that coordinate. This is outside the source's useful nonuniform-error parameter regime, not a counterexample to its reductions. At the other endpoint, known deterministic error is just an invertible public shift and adds no difficulty.
\subsection{Verification and remaining gaps}
On $\mathbb Z/7\mathbb Z$, point masses at 0 and 1 have distance 1. Deterministic noise leaves it 1, noise uniform on $\{0,1\}$ gives $1/2$, and uniform group noise gives 0; all convolutions were computed exactly. The failed inference is equating this statistical contraction with a strict comparison of computational hardness in the nondegenerate cryptographic regime. No efficient solver, strict hardness separation, or improvement of the source's parameterized reductions is produced.

\section{References and source attribution}
Primary workshop report: \url{https://doi.org/10.4171/owr/2014/35}. Relevant passages read: 2014/35 pp. 1952--1953, exact learning-with-errors-in-the-exponent samples, noise law and parameterized reductions.

All retrievals were checked on 11 October 2026. Original question attribution remains in the preserved Sources section. External theorems are attributed and not claimed as new discoveries or fully verified merely by citation.


\paragraph{Integration review.} Codex (GPT-6; exact serving revision not exposed) reviewed the written arguments and expanded layout and source-reference wrappers. Repository metadata and the production web renderer were checked using Python 3.13 and JavaScript with MathJax 3.2.2. This model-assisted review is not independent human proof verification; no PDF compilation or formal proof checking is asserted.

\section{Completion Estimate}
\noindent\textbf{COMPLETION: 0\%}
No substantiated new progress toward the entire original target is claimed. The worked results and failed route are delimited above; a known result, illustrative model, or auxiliary obstruction is not counted as original-target completion.

\section{Final status and attribution}
\textbf{Attempt status: unresolved.} This describes this research attempt, not a new ruling on the source catalogue's status or on external literature claims.
The mathematics and this self-audit were generated by OpenAI \texttt{gpt-6-astra}, reasoning \texttt{ultra}, on 11 October 2026. Codex (GPT-6; exact serving revision not exposed) orchestrated the batch, reviewed the written mathematical arguments, and contributed formatting/repository review. No human mathematical authorship, independent human proof checking, or formal verification is claimed. The historical preparation paragraphs in the reproduced catalogue body concern that body only.

Contribution prompt A was applied to the complete definition and supplied source record, followed by the adversarial checks of prompt B. Local reproduction files are \texttt{batch019/checks.py}, \texttt{computations.json}, and the per-ID audit JSON. They were not public links at preparation time. Numerical checks are finite and do not prove unrestricted results. No LaTeX engine was found locally; compilation was not performed by this research agent.

\end{document}
